\documentclass[../veristore_security_model.tex]{subfiles}

\begin{document}

\section{Scoping}
\subsection{In Scope}

This security model focuses on the core veri-store protocol and implementation. In scope are the parts of the system that directly handle data fragments, fingerprints, and cross-checksums:

\begin{itemize}
    \item Homomorphic fingerprint generation and verification for data blocks and fragments.
    \item Reed--Solomon erasure encoding and decoding (3-of-5 scheme) used to split and reconstruct data.
    \item Client logic that disperses data, generates the fingerprinted cross-checksum (fpcc), and verifies fragments on retrieval.
    \item Storage server API and logic for receiving, storing, verifying, and deleting fragments.
    \item The fingerprinted cross-checksum protocol, including how servers use it to check fragment integrity.
    \item Interactions between the veri-store client and storage servers over HTTP (FastAPI and \texttt{httpx}).
    \item Protection of data integrity and authenticity for stored fragments and reconstructed blocks.
\end{itemize}

\subsection{Out of Scope}

Some important security topics are intentionally out of scope for this project. We document them here so that readers do not assume they are handled:

\begin{itemize}
    \item Physical security of machines (e.g., data center access control, hardware tampering).
    \item Operating system and hypervisor security, including kernel exploits and container escapes.
    \item Network-layer protections such as DDoS mitigation, BGP attacks, or link encryption; we assume a working TCP/IP network and optionally TLS.
    \item User authentication and authorization beyond the ``authenticated client'' abstraction; account management and access control are not modeled.
    \item Confidentiality of stored data; veri-store detects tampering but does not encrypt data. Sensitive data must be encrypted before use.
    \item Availability in the face of more than two failed or unreachable servers (beyond the 3-of-5 erasure code tolerance).
    \item Advanced Byzantine behavior such as large-scale collusion between many servers or compromised clients coordinating attacks.
\end{itemize}

\newpage
\end{document}